Das Rätselspiel Tatamibari (japanisch
\begin{CJK}{UTF8}{min}タタミバリ\end{CJK}, von
\begin{CJK}{UTF8}{min}畳\end{CJK}
(\begin{CJK}{UTF8}{min}たたみ\end{CJK}, Tatami), japanische Reisstrohmatte
und \begin{CJK}{UTF8}{min}張り\end{CJK}
(\begin{CJK}{UTF8}{min}ばり\end{CJK}, bari)
mit der Bedeutung von ``im Stil von'')
wird auf einem $n\times m$-Spielfeld gespielt, in dem in einzelnen Feldern
die Symbole aus $\Sigma=\{\texttt{+}, \texttt{-}, \texttt{|}\}$ eingetragen
sind.
Der Spieler muss das Spielfeld entlang der Kanten zwischen den Feldern
in Rechtecke aufteilen und dabei die folgenden Regeln einhalten:
\begin{enumerate}
\item
In jedem Rechteck findet man genau ein Feld mit einem Symbol aus $\Sigma$.
\item
Enthält ein Rechteck ein \texttt{-}, ist es breiter als hoch.
\item
Enthält ein Rechteck ein \texttt{|}, ist es höher als breit.
\item
Enthält ein Rechteck ein \texttt{+}, ist es ein Quadrat.
\item
Vier Rechtecke dürfen nicht an einer gemeinsamen Ecke zusammentreffen.
\end{enumerate}
Die Abbildung zeigt ein Tatamibari-Rätsel (links) mit Lösung (rechts, die
Farben sind nur zur besseren Sichtbarkeit der Rechteck der Unterteilung).
\begin{center}
\def\punkt#1#2{({0.7*(#1)},{0.7*(#2)})}
\def\horizontal#1#2{
	\draw[line width=1.4pt] \punkt{#1+0.2}{#2+0.5} -- \punkt{#1+0.8}{#2+0.5};
}
\def\vertikal#1#2{
	\draw[line width=1.4pt] \punkt{#1+0.5}{#2+0.2} -- \punkt{#1+0.5}{#2+0.8};
}
\def\kreuz#1#2{
	\horizontal{#1}{#2}
	\vertikal{#1}{#2}
}
\def\feld{
	\draw \punkt{0}{0} rectangle \punkt{7}{7};
	\foreach \i in {1,...,6}{
		\draw \punkt{\i}{0} -- \punkt{\i}{7};
		\draw \punkt{0}{\i} -- \punkt{7}{\i};
	}
	\vertikal{1}{4}
	\kreuz{2}{1}
	\horizontal{2}{5}
	\vertikal{4}{1}
	\horizontal{4}{3}
	\horizontal{4}{6}
	\kreuz{5}{0}
	\vertikal{5}{5}
}
\definecolor{farbeeins}{rgb}{1.0,0.6,0.8}
\definecolor{farbezwei}{rgb}{0.8,0.6,1.0}
\definecolor{farbedrei}{rgb}{0.6,0.8,1.0}
\definecolor{farbevier}{rgb}{0.6,1.0,0.8}
\definecolor{farbefuenf}{rgb}{0.8,1.0,0.6}
\definecolor{farbesechs}{rgb}{1.0,0.8,0.6}
\definecolor{farbesieben}{rgb}{0.8,0.8,0.8}
\definecolor{farbeacht}{rgb}{1.0,1.0,0.8}
\begin{tikzpicture}[>=latex,thick]
\begin{scope}[xshift=-3cm]
\feld
\end{scope}
\begin{scope}[xshift=3cm]
\fill[color=farbeeins!40] \punkt{0}{0} rectangle \punkt{3}{3};
\fill[color=farbezwei!40] \punkt{3}{0} rectangle \punkt{5}{3};
\fill[color=farbedrei!40] \punkt{5}{0} rectangle \punkt{7}{2};
\fill[color=farbevier!40] \punkt{5}{2} rectangle \punkt{7}{6};
\fill[color=farbefuenf!40] \punkt{0}{3} rectangle \punkt{2}{7};
\fill[color=farbesechs!40] \punkt{2}{3} rectangle \punkt{5}{5};
\fill[color=farbesieben!40] \punkt{2}{5} rectangle \punkt{5}{6};
\fill[color=farbeacht] \punkt{2}{6} rectangle \punkt{7}{7};
\feld
\draw[line width=2pt] \punkt{2}{3} -- \punkt{2}{7};
\draw[line width=2pt] \punkt{3}{0} -- \punkt{3}{3};
\draw[line width=2pt] \punkt{5}{0} -- \punkt{5}{6};
\draw[line width=2pt] \punkt{5}{2} -- \punkt{7}{2};
\draw[line width=2pt] \punkt{0}{3} -- \punkt{5}{3};
\draw[line width=2pt] \punkt{2}{5} -- \punkt{5}{5};
\draw[line width=2pt] \punkt{2}{6} -- \punkt{7}{6};
\end{scope}
\end{tikzpicture}
\end{center}
Kann eine nichtdeterministische Turing-Maschine in polynomieller
Zeit entscheiden, ob ein Tatamibari-Rätsel eine Lösung hat?

\begin{loesung}
Die Frage ist gleichbedeutend damit, ob Tatamibari in NP liegt.

In einem Tatamibari-Spielfeld gibt es $N=(n-1)m+n(m-1)=2nm-n-m$
Liniensegmente, aus denen man Unterteilungen erzeugen kann.
Indem man alle $2^N$ möglichen Auswahlen von Liniensegmenten
durchprobiert und auf die Einhaltung der genannten Regeln prüft.

Tatamibari liegt genau dann in NP, wenn es einen polynomiellen
Verifizierer dafür gibt.
Als Lösungszertifikat wird die Unterteilung des Spielfelds
verlangt.
Zusätzlich sollen alle Felder eines Teils mit der gleichen Farbe
einfärbt sein, die jedoch für kein anderes Teil des Feldes verwendet
wird.
Damit wird dann die folgende Verifikation durchgeführt.
Wir nennen dabei ein Feld berandet auf einer Seite berandet, wenn
entweder die Kante Teil einer Berandungslinie oder des Rands
des Spielfelds ist.
\begin{center}
\begin{tabular}{|c|p{11cm}|>{$}c<{$}|}
\hline
Schritt&Verifikation&\text{Aufwand}\\
\hline
1&
Für jedes Feld, welches links und oben berandet ist, folge den von der linken
oberen Ecke ausgehenden Randsegmenten nach unten, wähle jeweils Abzweigungen
nach links wenn vorhanden, niemals aber nach rechts, und überprüfe, ob 
die vierte Richtungsänderung im ursprünglichen Eckpunkt erfolgt.
&O(n^2m^2)\\
2&
Für jedes Feld, welches ein \texttt{-} enthält, zähle die
horizontal und vertikal benachbarten Felder gleicher Farbe und
überprüfe, dass die horizontale Anzahl grösser ist als die vertikale.
&O(nm)\\
3&
Für jedes Feld, welches ein \texttt{|} enthält, zähle die
horizontal und vertikal benachbarten Felder gleicher Farbe und
überprüfe, dass die vertikale Anzahl grösser ist als die horizontale.
&O(nm)\\
4&
Für jedes Feld, welches ein \texttt{+} enthält, zähle die
horizontal und vertikal benachbarten Felder gleicher Farbe und
überprüfe, dass beide Anzahlen gleich sind.
&O(nm)\\
5&
Für jede Ecken kontrolliere, dass die an dieser Ecke zustammenstossenden
Felder nicht vier verschiedene Farben haben.
&O(nm)\\
\hline
&Total&O(n^2m^2)\\
\hline
\end{tabular}
\end{center}
Offenbar ist die Verifikation in polynomieller Zeit ausführbar, es gibt
also einen polynomiellen Verifizierer und damit ist gezeigt,
dass Tatamibari in NP ist.
\end{loesung}

\begin{diskussion}
Tatamibari ist NP-vollständig: \url{https://arxiv.org/pdf/2003.08331}.
\end{diskussion}


\begin{bewertung}
Entscheidbarkeit ({\bf E}) 1 Punkt,
Prinzip Verifizierer ({\bf V}) 1 Punkt,
Zertifikat spezifiziert ({\bf Z}) 1 Punkt,
Laufzeitschätzung ({\bf L}) 2 Punkte,
Schlussfolgerung polynomieller Verifizierer ({\bf S}) 1 Punkt.
\end{bewertung}

